\begin{table*}[t]
\centering\footnotesize
\setlength{\tabcolsep}{3pt}
\caption{최신 사람 가시성 상태에 따른 YOLO Base와 N3 보조 분석. 원래2499개 참조 좌표를 유지한다.}\label{sup:visibility}
\begin{tabularx}{\textwidth}{Ylrrrr}
\toprule
상태 & 방법 & \shortstack{참조/유효\\예측 코너} & 중앙값(px) & P90(px) & \pckten(\%) \\
\midrule
직접 가시 & Base & 1,776/1,759 & 5.829 & 25.761 & 70.946 \\
직접 가시 & N3 & 1,776/1,759 & 4.975 & 23.571 & 77.646 \\
외부 가림 & Base & 218/195 & 19.485 & 80.896 & 27.523 \\
외부 가림 & N3 & 218/195 & 18.269 & 81.568 & 25.688 \\
자체 가림 & Base & 462/448 & 8.874 & 41.132 & 53.247 \\
자체 가림 & N3 & 462/448 & 8.143 & 38.356 & 55.988 \\
화면 밖 & Base & 43/43 & 12.159 & 224.325 & 44.186 \\
화면 밖 & N3 & 43/43 & 10.864 & 223.509 & 47.287 \\
\bottomrule
\end{tabularx}
\tabnote{직접 가시1776·외부 가림218·자체 가림462·화면 밖43점으로 합계2499점이다. 결측·객체 매칭 실패는 전체 참조 PCK 분모에 유지했다. 조건부 중앙값·P90와 전체 분모를 구분한다. 기존 잠금 상태71점은 덮어쓰지 않았고, 가시성 상태를 입력해도 원래 참조 생성 경로가 바뀌지는 않는다. P·N2 및 DOPE·ResNet·정사각형 두 모드의 모든 행은 별도CSV에 보존한다.}
\end{table*}
